% Comentarios interpretativos, escritos a mano después de ver los datos. Los números vienen de las macros
% \R{..}{..} generadas por `python -m unirl tables`; aquí no se copia ninguno.

\newcommand{\RespuestaAdelgazar}{Se comprobó con \texttt{tools/sh\_goal.py}: con una sola premisa, y tanto sin como con los tres hechos previos de la demostración, Sledgehammer tampoco la encuentra en \R{SH}{timeout}\,s (encontrada: \R{SH}{sin_preludethin}, sin los hechos; \R{SH}{con_prelude_del_paperthin}, con ellos). En esta máquina, adelgazar las premisas no basta: la línea de los autores lleva una selección de 18 hechos que la automatización no reproduce.}

\newcommand{\BaseIComentario}{Con el invariante y el variante del artículo, el juez $V_I$ de 09 da \lab{\R{BaseI}{label}} con \R{BaseI}{closed} de \R{BaseI}{nvcs} VCs cerradas: quedan abiertas las \R{BaseI}{openvar} del variante y las \R{BaseI}{openpost} de la postcondición. Es el mismo recuento que en 08 sin \id{bisect_root_extract}, así que los tres hechos previos no cambian nada. Este es el listón para el agente: con la anotación de los autores, el juez tampoco acepta.}

\newcommand{\ComentarioGrafo}{Bajo $V_G$, \texttt{gpt-6-luna} acepta todos los episodios de los cinco métodos (20 de 20 en cada uno; Tabla~\ref{tab:resumen}), casi siempre al primer intento: $t^*$ mediano 1 en los cinco. Los pocos fallos son casi todos de orden (\lab{Invalid\_Step}): el modelo usa un paso antes de establecer sus dependencias y corrige con la retroalimentación. Sólo en secante cae una vez en un distractor (\R{SecGLuna}{DISTRACTOR} en \R{SecGLuna}{att} intentos). Con el texto de los pasos a la vista y la longitud mínima anunciada, $V_G$ mide sobre todo si el modelo respeta dependencias explícitas, no si sabe demostrar.}

\newcommand{\ComentarioComparacion}{Con el mismo proponente y el mismo presupuesto, los dos jueces dan resultados opuestos. Bajo $V_G$, \texttt{gpt-6-luna} acepta \R{BisGLuna}{acc} de \R{BisGLuna}{ep} episodios con $t^*$ mediano \R{BisGLuna}{tstar}; bajo $V_I$, \R{BisILuna}{acc} de \R{BisILuna}{ep}, con sus \R{BisILuna}{nprops} intentos en \lab{VC\_Open}. La mejor propuesta cerró \R{BisILuna}{bestcn} VCs y la media fue \R{BisILuna}{meanfrac}, frente a \R{BaseI}{closed}/\R{BaseI}{nvcs} con la anotación del artículo. Cada intento de $V_I$ costó en promedio \R{BisILuna}{judges}\,s de juez, frente a \R{BisILuna}{llms}\,s de modelo.

Tres rasgos de las propuestas explican dónde se atasca. Primero, \R{BisILuna}{nexists} de \R{BisILuna}{nprops} meten la existencia de la raíz en el invariante ($\exists c \in [l,u].\ f(c) = 0$), como en una demostración de libro. Eso traslada el argumento de Bolzano a la preservación del invariante, y en las \R{BisILuna}{nexistsopen} esa VC queda abierta: Sledgehammer no la cierra en 60\,s. El artículo hace lo contrario: deja la raíz fuera del invariante y la extrae al final con un lema escrito a mano. Segundo, \R{BisILuna}{nulvar} usan un variante en función del intervalo, $\lceil \log_2((u-l)/\mathit{tol}) \rceil$, y \R{BisILuna}{nitervar} uno en función de \texttt{iter}, como el del artículo; en ambos casos las VCs del variante que quedan abiertas (\R{BisILuna}{nvaropenlog} de \R{BisILuna}{nvaropen}) son aritmética con $\lceil \log_2 \cdot \rceil$, la misma que los autores resuelven con su línea \texttt{smt} de 18 hechos. Tercero, en los \R{BisILuna}{npostopen} intentos con VCs quedó abierta la postcondición: sin \id{bisect_root_extract}, ninguna anotación consigue que la automatización extraiga la raíz. Sumando sobre los intentos, las VCs abiertas fueron \R{BisILuna}{openinv} de invariante, \R{BisILuna}{openvar} de variante y \R{BisILuna}{openpost} de postcondición (clasificación heurística por la conclusión y las premisas de cada VC). Ningún invariante coincide con el del artículo (\R{BisILuna}{nsameinv} de \R{BisILuna}{nprops}); el variante sí, literalmente, en \R{BisILuna}{nsamevar}, y es la elección natural (Apéndice~\ref{app:propuestas}).

La diferencia entre jueces no dice que $V_I$ sea ``más difícil'' en abstracto. Dice dónde está el trabajo: en $V_G$ el profesor ya eligió los lemas, cerró la aritmética y escribió el argumento de Bolzano, y el agente sólo los ordena; en $V_I$ esas tres cosas siguen abiertas, y la automatización de una línea no las hace sola ni con la anotación de los autores. Para que $V_I$ sea un entorno de aprendizaje útil hace falta darle al agente más acciones que la anotación: selección de premisas, lemas auxiliares propuestos por él y verificados, o pruebas Isar (experimento 5).}

